\section{Zusammenfassung}
Zusammenfassend war diese Diplomarbeit ein lehrreiches Projekt, bei dem viele neue
Erfahrungen gesammelt wurden. Gleichzeitig befasst sich diese Diplomarbeit mit der Konzeption,
Implementierung und Dokumentation des Systems \textbf{DropIT}, einer webbasierten
Lösung zur strukturierten, sicheren und effizienten Verwaltung digitaler
Dokumente. Ziel der Arbeit war es, die Ablage, Klassifikation, Suche und
Organisation von Dokumenten in einer einheitlichen Anwendung zu bündeln und
dabei sowohl technische als auch organisatorische Anforderungen moderner
Dokumentenverwaltung zu berücksichtigen.

Im Rahmen der Arbeit wurde eine serviceorientierte Systemarchitektur umgesetzt,
die Benutzeroberfläche, Fachlogik, Identitätsverwaltung und externe Dienste
klar voneinander trennt. Die Webanwendung basiert auf Blazor und wird über
Microsoft Entra ID authentifiziert. Für die eigentliche Dokumentablage wird
SharePoint Online verwendet, während strukturierte Metadaten in PostgreSQL
gespeichert und Suchanfragen über Azure AI Search verarbeitet werden. Die
zentrale Orchestrierung der Geschäftslogik erfolgt über den Zone-Service, der
den Upload-Prozess, die Metadatenverarbeitung sowie die Anbindung der weiteren
Backend-Dienste steuert.

Ein wesentlicher Funktionsbaustein von DropIT ist die KI-gestützte
Dokumentenklassifikation. Über den AI-Service werden hochgeladene PDF-Dokumente
analysiert, automatisch einem Dokumenttyp zugeordnet und um relevante Metadaten
wie Jahr, Vertragstyp, Partner oder Ablaufdatum ergänzt. Dadurch wird der
manuelle Erfassungsaufwand reduziert und eine konsistentere Weiterverarbeitung
der Dokumente ermöglicht. Ergänzend dazu wurde eine Sucharchitektur realisiert,
die Volltextsuche, Metadatenfilter und benutzerbezogene Zugriffskontrolle
kombiniert, um Dokumente schnell und zielgerichtet auffindbar zu machen.

Die Arbeit zeigt damit, wie sich Microsoft-Dienste, moderne Webtechnologien und
KI-gestützte Dokumentenanalyse zu einer praxistauglichen Gesamtlösung verbinden
lassen. DropIT stellt eine belastbare Grundlage für eine zentrale
Dokumentenverwaltung dar und verdeutlicht, wie Automatisierung, Suchbarkeit und
Sicherheit in einem gemeinsamen System zusammengeführt werden können.
